\begin{claim}{F.26}{Proved}{The exact trace obstruction and all affine marginal images with scalar field of order nine}
Retain the actual inclusions
$\mathcal K_\pm\leq R_\pm\leq V_\pm$ of F.19--F.25.
Fix either sign, and temporarily write these groups as
$K\leq R\leq V$. A proper affine marginal image
$\rho(K)\subsetneq\rho(R)$, for a continuous
$\rho:V\to\operatorname{Aff}_1(k)$ over a finite field,
has an actual continuous affine image of the following form:
\[
 \begin{gathered}
 \rho':V\longrightarrow k_0\rtimes D,\qquad
 \rho'(V)=k_0\rtimes D,\qquad \rho'(R)=k_0\rtimes C,\\
 \rho'(K)\ \leq\
       \ker\bigl(\operatorname{Tr}_{k_0/k_1}:k_0\to k_1\bigr)
                       \rtimes C,\\
 C=\chi(R)=\chi(K)\ne\{1\},\qquad
 D=\chi(V),\qquad
 k_1=\mathbb F_p(C)\subsetneq k_0=\mathbb F_p(D).
 \end{gathered}                                                \tag{F.26.1}
\]
Here $\chi$ is the original scalar character, retained
by the explicit conjugation and coefficient map constructed
below. In particular the last group in (F.26.1) is a
proper subgroup of the displayed $R$ image.

For any finite-field data $k_0,\chi,C,k_1$ as in
(F.26.1), including the equality case $k_1=k_0$, put
$q_1=|k_1|$ and $m=[k_0:k_1]$. The exact coefficient map is
\[
 \begin{aligned}
 \mathscr T_\chi:
 Z^1_{\rm cont}(V,k_{0,\chi})&\longrightarrow
 Z^1_{\rm cont}(K,k_{1,\chi|K}),\\
 d&\longmapsto\operatorname{Tr}_{k_0/k_1}(d|K).
 \end{aligned}                                                \tag{F.26.2}
\]
Its known kernel subspace is
\[
 \mathscr P_\chi=
 \{\,g\mapsto(\chi(g)-1)b:
                      b\in k_0,\ \operatorname{Tr}(b)=0\,\}.
                                                               \tag{F.26.3}
\]
Every element of $\ker\mathscr T_\chi\setminus\mathscr P_\chi$
gives a proper actual affine marginal image, and every
proper affine marginal image gives such an element for
the data in (F.26.1). Thus
$\ker\mathscr T_\chi=\mathscr P_\chi$ for all the stated
data is an exact remaining criterion for the full
scalar affine image equality. The proof gives an explicit
coefficient isomorphism identifying this criterion with
the summed restrictions for the distinct characters
$\chi^{q_1^j}$, $0\leq j<m$. Its rank target is $6m+1$
on $7m$ columns; no general attainment of that rank
is asserted here.

For the scalar field of order nine that rank is attained.
More precisely, for every finite field $k$ and every
continuous $\rho:V_\pm\to\operatorname{Aff}_1(k)$ whose
scalar values generate a subfield of order nine,
\[
                         \rho(\mathcal K_\pm)=\rho(R_\pm).
                                                               \tag{F.26.4}
\]
Consequently this equality holds for every continuous
map $V_\pm\to\operatorname{Aff}_1(\mathbb F_9)$, including
maps whose scalar field is $\mathbb F_3$. The finite
calculation proving the additional field case exhausts
exactly $7\,560$ character tuples and retains the
literal F.23 source automorphisms.
\end{claim}

\begin{proof}
\emph{Reduction through the actual affine image.}
Write $H=\rho(V)$, $N=\rho(R)$, $M=\rho(K)$ and
let $D,C$ be their scalar images as above. The F.21
normal generation, transported through the actual
$D^+\to R$ of F.22, gives
\[
 \langle\!\langle M\rangle\!\rangle_N=N,\qquad
                         \chi(K)=\chi(R)=C.                   \tag{F.26.5}
\]
These are exactly the maps and equalities proved
in F.25.5. A proper marginal requires $C\ne\{1\}$,
since $C=\{1\}$ makes $N$ abelian. Put
$k_0=\mathbb F_p(D)$ and $k_1=\mathbb F_p(C)$.
F.25 excludes equality of these two fields for
a proper marginal.

Choose a generator $c$ of $C$ and an actual lift
$(b,c)\in M$. Set $z=b/(1-c)$ and conjugate the
representation by $\tau_z^{-1}$, where
$\tau_z(w)=w+z$. As in F.25.7--8, this gives
\[
 M_z=U\rtimes C,\qquad
 U=\{u:(u,1)\in M\},\qquad (0,c)\in M_z,
                                                               \tag{F.26.6}
\]
with $U$ unchanged and $k_1$-linear. Let $A_H$ and
$A_N$ be the translation subgroups of $H_z$ and
$N_z$. Because $N_z$ is normal in $H_z$, conjugation
by every actual scalar lift in $H_z$ preserves both
subgroups and multiplies translations by its scalar.
Thus both are vector subspaces over
$\mathbb F_p[D]=k_0$. For $a\in A_H$ the exact commutator is
\[
                  [(a,1),(0,c)]=((1-c)a,1)\in N_z .
                                                               \tag{F.26.7}
\]
Since $1-c\ne0$ and $A_N$ is $k_0$-linear, this proves
$A_H=A_N=:A$. For an arbitrary $(b',d)\in H_z$ the
same multiplication gives
$[(b',d),(0,c)]=((1-c)b',1)\in A$ and hence $b'\in A$.
Multiplying by its inverse translation then gives
the corresponding pure scalar element. Consequently
\[
 H_z=A\rtimes D,\qquad N_z=A\rtimes C,\qquad
 M_z=U\rtimes C,\qquad U\subsetneq A .                          \tag{F.26.8}
\]
Here $A$ is $k_0$-linear and $U$ is $k_1$-linear.

Choose a $k_1$-linear functional $L:k\to k_1$
vanishing on $U$ and nonzero on $A$. To construct
one, extend a basis of $U$ by an element of
$A\setminus U$, then to a basis of $k$, and give
that selected element value one and the other
basis vectors value zero.

\emph{The precise trace and functional factorization.}
Write $m=[k_0:k_1]$ and $q_1=|k_1|$. The trace is
\[
 \operatorname{Tr}(x)=\sum_{j=0}^{m-1}x^{q_1^j}.
                                                               \tag{F.26.9}
\]
Its $q_1$th power equals itself, using
$x^{q_1^m}=x$, so its values lie in $k_1$.
It is $k_1$-linear and nonzero: its polynomial
has degree $q_1^{m-1}<|k_0|$ and leading coefficient
one, and therefore cannot vanish on all of $k_0$.
It is consequently onto $k_1$. Choose $\beta$
with $\operatorname{Tr}(\beta)=1$. For every
$a\ne0$,
$\operatorname{Tr}(a(\beta/a))=1$. Thus the pairing
$\operatorname{Tr}(ab)$ is nondegenerate, and
\[
 k_0\xrightarrow{\ \sim\ }
 \operatorname{Hom}_{k_1}(k_0,k_1),\qquad
 a\longmapsto(b\longmapsto\operatorname{Tr}(ab))
                                                               \tag{F.26.10}
\]
is an isomorphism: it is injective and both sides
have $k_1$-dimension $m$. This also covers characteristics
dividing $m$.

Choose a $k_0$-basis $v_1,\ldots,v_r$ of $k$.
For each $i$, (F.26.10) gives a unique $a_i\in k_0$
with $L(tv_i)=\operatorname{Tr}(a_it)$ for all $t\in k_0$.
Define
\[
 \ell:k\longrightarrow k_0,\qquad
 \ell\left(\sum_i t_iv_i\right)=\sum_i a_it_i,\qquad
                         L=\operatorname{Tr}\ell .
                                                               \tag{F.26.11}
\]
This is the required $k_0$-linear coefficient map.
Its image on $A$ is a nonzero $k_0$-subspace of
$k_0$, hence is all of $k_0$; its image on $U$
lies in $\ker\operatorname{Tr}$. The exact homomorphism
of finite affine groups is
\[
 P_\ell:k\rtimes D\longrightarrow k_0\rtimes D,\qquad
                         (b,d)\longmapsto(\ell(b),d).
                                                               \tag{F.26.12}
\]
The identity $\ell(db)=d\ell(b)$ proves its compatibility
with multiplication. All groups here are finite and
discrete, so it is continuous. Composing it with the
specified conjugation and with $\rho$ gives $\rho'$.
Equations (F.26.8)--(F.26.12) prove every image assertion
in (F.26.1), with the original scalar character retained.

\emph{The exact trace-cocycle criterion.}
For the data in (F.26.2), the restriction of $\chi$
to $K$ lies in $C\subset k_1^\times$. Thus composing
a crossed homomorphism with the $k_1$-linear trace
on $K$ proves precisely the target cocycle equation
in (F.26.2); continuity is preserved. Each principal
cocycle in (F.26.3) maps to zero. Because $C\ne\{1\}$,
the parameter map $b\mapsto(\chi-1)b$ is injective,
and this subspace has $k_1$-dimension $m-1$.

The representation $\rho'$ from (F.26.1) has a
translation cocycle $d'$ in $\ker\mathscr T_\chi$.
It is not principal: a principal affine representation
fixes a point of the field, whereas $\rho'(V)$ contains
every translation, including nonzero translations
which move every point. This gives an element outside
$\mathscr P_\chi$.

Conversely take $d\in\ker\mathscr T_\chi
\setminus\mathscr P_\chi$. The specified homomorphism
\[
 \rho_d:V\longrightarrow k_0\rtimes D,\qquad
                        g\longmapsto(d(g),\chi(g))             \tag{F.26.13}
\]
is continuous by the cocycle rule. Its translation
subgroup is stable under $D$ and hence is a
$k_0$-subspace of $k_0$. It cannot be zero. Otherwise
the scalar projection would identify $\rho_d(V)$
with the cyclic group $D$. A generator lift $(b,a)$,
$a\ne1$, fixes $b/(1-a)$, and all its powers fix that
same point. Then $d$ is principal; its vanishing
trace on $K$, together with $C\ne1$, makes the
parameter trace zero, contrary to the choice of $d$.
Thus $\rho_d(V)$ contains all translations of $k_0$.
Its $R$ image is normal and has nontrivial scalar
image $C$. Formula (F.26.7) therefore forces all
translations into that image, giving
$\rho_d(R)=k_0\rtimes C$. Its $K$ image lies in
$\ker\operatorname{Tr}\rtimes C$, by (F.26.2), and
is proper. This proves both directions of the criterion.

\emph{The coefficient splitting, with its inverse.}
The characters $\chi_j=\chi^{q_1^j}$, $0\leq j<m$,
are pairwise distinct: a nonidentity Frobenius power
fixing all their scalar values would fix their
generated field $k_0$. They all agree on $R$.
The exact coefficient isomorphism is
\[
 S:k_0\otimes_{k_1}k_0\xrightarrow{\ \sim\ }\prod_{j=0}^{m-1}k_0,
             \qquad b\otimes a\longmapsto
                       (a b^{q_1^j})_{j=0}^{m-1}.              \tag{F.26.14}
\]
Choose a multiplicative generator $\gamma$ of $k_0^\times$.
Its minimal polynomial $f$ over $k_1$ has degree $m$
and distinct roots $\gamma^{q_1^j}$. Identify the
tensor ring with $k_0[X]/f(X)$, with the first
$\gamma$ represented by $X$ and the second factor
supplying coefficients. Evaluation at these roots
is (F.26.14), and its inverse is
\[
 (a_j)_j\longmapsto
 \sum_{j=0}^{m-1} a_j
   \prod_{r\ne j}
     \frac{X-\gamma^{q_1^r}}{\gamma^{q_1^j}-\gamma^{q_1^r}}
                      \pmod{f(X)}.                            \tag{F.26.15}
\]
Every denominator is nonzero and the two composites
are the identity by evaluation and degree.

Give the first tensor factor the $\chi$ action.
Under $S$ its $j$th factor has the $\chi_j$ action.
A continuous cocycle on $V=\widehat F_7$ is specified
by its seven values: any such tuple defines a map
from the discrete free group into the finite affine
group and extends uniquely to its full completion.
Thus (F.26.14) gives the exact comparison
\[
 \begin{aligned}
 Z^1_{\rm cont}(V,k_{0,\chi})\otimes_{k_1}k_0
    &\xrightarrow{\ \sim\ }
       \bigoplus_{j=0}^{m-1}Z^1_{\rm cont}(V,k_{0,\chi_j}),\\
 d\otimes a&\longmapsto(a d^{q_1^j})_j .
 \end{aligned}                                                \tag{F.26.16}
\]
Expansion in a $k_1$-basis of $k_0$ likewise identifies
the scalar extension of the target of (F.26.2)
with $Z^1_{\rm cont}(K,k_{0,\chi|K})$. Coordinate
projections preserve cocycles and continuity, so
this target comparison holds for the actual
profinite subgroup $K$ as well. On simple tensors,
the trace formula proves that $\mathscr T_\chi\otimes1$
is exactly
\[
                  (d_j)_j\longmapsto\sum_j d_j|K.              \tag{F.26.17}
\]
The scalar extension of $\mathscr P_\chi$ is
$\{(b_j(\chi_j-1))_j:\sum_jb_j=0\}$, since
the trace becomes coordinate summation under $S$.
Extension by a field preserves kernels and detects
a nonzero vector space by a basis. This proves
the claimed equivalent kernel description and
the rank target $7m-(m-1)=6m+1$.

\emph{The exhaustive field-of-order-nine character set.}
Use the actual negative coordinates of F.24.4 and
write $\mathbb F_9=\mathbb F_3[\iota]/(\iota^2+1)$.
Every element is uniquely $a+b\iota$ with
$a,b\in\mathbb F_3$; the polynomial has no root
in $\mathbb F_3$, so this quotient is a field.
For any scalar field $k_0$ of order nine, choose
an element $\iota_0$ of order four in $k_0^\times$.
It satisfies $\iota_0^2=-1$. The explicit field map
$a+b\iota\mapsto a+b\iota_0$ is an isomorphism:
its source is a field, it preserves one, and both
fields have nine elements. Its inverse takes the
unique two coefficients in this basis. Apply that
inverse to both affine coordinates to obtain the
corresponding isomorphism of affine groups. It
preserves the original source representation and
its two subgroup images through this specified
coefficient map. Either choice of $\iota_0$ is
covered by the complete character enumeration.
The field code $a+3b$ in the programs records these
two coordinates exactly. Multiplication is
\[
 (a+b\iota)(c+d\iota)=(ac-bd)+(ad+bc)\iota,
 \qquad
 \operatorname{Tr}(a+b\iota)=2a,\qquad
 \operatorname{Tr}(\iota(a+b\iota))=b .
                                                               \tag{F.26.18}
\]
Thus multiplication of cocycles by $\iota$ is the
exact coefficient isomorphism from the imaginary-part
kernel computed below to the trace kernel, with
inverse multiplication by $-\iota$. It sends
real principal parameters to the trace-zero
principal parameters $\mathbb F_3\iota$.

F.25 covers trivial $C$ and equality of scalar
fields. The remaining scalar-field-nine data
have $k_1=\mathbb F_3$, so $C=\{1,-1\}$.
The four deleted character values are in
$\{1,-1\}$ with at least one equal to $-1$.
Conversely those values give exactly this $C$,
because the four deleted generators normally
generate $R$ and the character is abelian.
The three retained values lie in
$(\mathbb F_9^\times)^3$ and are not all in
$(\mathbb F_3^\times)^3$. These conditions give
$k_0=\mathbb F_9$, because the extension of
$\mathbb F_3$ has prime degree two. The exhaustive
character count is therefore
\[
                   (8^3-2^3)(2^4-1)=7\,560 .                  \tag{F.26.19}
\]
The first producer enumerates every such tuple.

For a fixed tuple $t$, write a cocycle-value vector
as $v=x+\iota y$, with $x,y\in\mathbb F_3^7$,
and put $\widetilde v=(x,y)^{\mathsf T}$.
For each of the sixty-four actual F.24 source rows
$a(t)=a_{\rm r}+\iota a_{\rm i}$ the exact condition
that its value have imaginary part zero is
\[
       (a_{\rm i},a_{\rm r})\widetilde v=0 .                   \tag{F.26.20}
\]
The seven-coordinate vector $t-\mathbf1$ with real
principal parameter gives a nonzero solution.
It is nonzero because some deleted value equals
$-1$. Every row arises from an actual element
of $K$, whose character lies in $\mathbb F_3$,
so this principal line is contained in the
kernel of all the rows and of every further
source conjugate.

The first exact row reduction gives the following
complete distribution:
\[
 \begin{array}{c|rrr}
 \text{rank on fourteen columns}&11&12&13\\\hline
 \text{character tuples}&8&396&7156 .
 \end{array}                                                  \tag{F.26.21}
\]
For the $7\,156$ rank-thirteen tuples, the kernel
is precisely the stated principal line.
The other $404$ tuples are retained explicitly
and passed to the second producer.

\emph{All eleven actual source actions on the remaining tuples.}
For each of these $404$ tuples, evaluate all eleven
matrices $A_i$ from F.24.9 by the same exact
Laurent evaluation. Write $A_i=A_{i,\rm r}
+\iota A_{i,\rm i}$. Its actual action on the
fourteen prime-field coordinates is
\[
 \widetilde A_i=
 \begin{pmatrix}
 A_{i,\rm r}&-A_{i,\rm i}\\
 A_{i,\rm i}& A_{i,\rm r}
 \end{pmatrix}.                                               \tag{F.26.22}
\]
Start with the eight original negative seed
imaginary-part rows. At each round replace their
row space $W$ by
\[
                 W+\sum_{i=0}^{10}W\widetilde A_i .
                                                               \tag{F.26.23}
\]
Every row lies in the $\mathbb F_3$-linear span of the
imaginary-part evaluation rows of actual source
conjugates of the seeds.
The pure source action preserves the character,
and its action on cocycles is exactly the
precomposition matrix proved in F.24.9.
The successive exact ranks are
\[
 \begin{array}{c|rrrr}
 \text{stage}&\multicolumn{4}{c}{\text{rank : number of tuples}}\\\hline
 \text{eight seeds}&5:80&6:284&7:2&8:38\\
 \text{one round}&12:50&13:354&&\\
 \text{two rounds}&13:404&&&\\
 \text{third round, stable}&13:404&&&
 \end{array}                                                  \tag{F.26.24}
\]
In particular two rounds already supply rank
thirteen for every remaining tuple. The third
round verifies invariance under all eleven
matrices. Their actions are invertible, since
they come from the actual source automorphisms.
An invariant subspace for an invertible linear
map on a finite-dimensional space is also
invariant under its inverse. Hence this stable
row space accounts for both directions of all
source generators.

All arithmetic here is in the explicitly described
finite field or in $\mathbb F_3$. Every nonzero
Laurent variable is evaluated using its powers
modulo eight. The field table, all $7\,560$
character choices, all $64$ row evaluations,
the $404$ retained tuples, the eleven actual
source matrices and the complete rank histories
are specified by the accompanying producers
and their fixed input files. Gaussian elimination
uses the exact inverses one and two in
$\mathbb F_3$. The independent verification imports
neither producer. It reconstructs these matrices
from the fixed universal source rows and actions,
using a separate multiplicative-log Laurent evaluation,
and computes the ranks by scalar incremental elimination
over $\mathbb F_3$. It verifies the entire character set,
all principal vectors, all eleven source actions,
every closure image, the stable row spaces and all
nine coefficient identities in (F.26.18). The universal
source data themselves are the audited literal source
evaluations of F.24.

It follows that every continuous cocycle whose
imaginary part vanishes on $K$ has values
$b(t-\mathbf1)$ for a unique $b\in\mathbb F_3$.
Indeed all the actual rows (F.26.20) and their
source conjugates must vanish on it, so the
proved ranks give precisely that one-dimensional
kernel. Conversely those principal cocycles
do have imaginary part zero on $K$.
Multiplication by $\iota$ in (F.26.18), and its
specified inverse, prove
$\ker\mathscr T_\chi=\mathscr P_\chi$ for all
$7\,560$ remaining field-of-order-nine characters.

The criterion proved in (F.26.1)--(F.26.13)
now excludes every proper marginal whose
scalar field has order nine, even when its
original coefficient field $k$ is larger.
For a representation into
$\operatorname{Aff}_1(\mathbb F_9)$ its scalar
field is either $\mathbb F_3$ or $\mathbb F_9$;
F.25 handles the former and the present argument
handles the latter. Finally, precomposition
with the actual continuous isomorphism
$\mathcal B_+:V_-\to V_+$ of F.23.13, and
its inverse, transports each positive
representation to the proved negative case,
carrying both $K$ and $R$ exactly as stated
there. This proves (F.26.4) for both signs.
\end{proof}
